% 《向量》中文版 PDF 的 LaTeX 导言区（由 tools/build_pdf.py 用 -H 注入）
% 每行配置都对应一个实测结论，改动前请看注释。

% ============ 数学 ============
\usepackage{amsmath,amssymb}
\usepackage{unicode-math}
% fontspec 认不出 "Latin Modern Math" 这个字体名（实测报 fontspec Error），
% 必须用 TeX Live 里的文件名。
% Scale 与下面的西文缩放保持一致：数学字体不受 \setmainfont 影响，不跟着缩的话
% 行内 $x$ 会比周围的西文大一圈（公式里的字母也是「西文」）。
\setmathfont{latinmodern-math.otf}[Scale=0.94]

% ============ 西文：略小于中文 ============
% 同一名义字号下，拉丁字母的视觉体量大于汉字（实测：11pt 时拉丁墨迹高
% 10.91pt、汉字 10.54pt；再加粗体/斜体的笔画对比，混排时西文会「跳出来」）。
% 中文出版惯例把西文缩到中文的 90–95%，这里取 0.94。
% 用文件名而不是字体名——与上面 \setmathfont 同一个原因。
\setmainfont{lmroman10-regular.otf}[
  Scale=0.94,
  BoldFont=lmroman10-bold.otf,
  ItalicFont=lmroman10-italic.otf,
  BoldItalicFont=lmroman10-bolditalic.otf]
\setmonofont{lmmono10-regular.otf}[Scale=0.94]
\setsansfont{lmsans10-regular.otf}[Scale=0.94]

% ============ 版式 ============
\usepackage{geometry}
\geometry{
  a4paper,
  inner=3.3cm, outer=2.2cm, top=2.5cm, bottom=2.5cm,
  marginparsep=3mm, marginparwidth=1.1cm
}
\usepackage{fancyhdr}
\pagestyle{fancy}
\fancyhf{}
\fancyhead[LE,RO]{\thepage}
\fancyhead[LO]{\nouppercase{\rightmark}}
\fancyhead[RE]{\nouppercase{\leftmark}}
\renewcommand{\headrulewidth}{0.4pt}
% 章号由 build_pdf.py 写进 H1，LaTeX 不再自动编号；
% 不能用 \CTEXthechapter —— ctex 的 \chapter 在自增计数器之前就调用 \chaptermark，
% 取到的永远是「第零章」。
% pandoc 会设 secnumdepth=-\maxdimen，此时 \chapter 是否步进 chapter 计数器
% 取决于实现，导致「按章重置图表/脚注」时灵时不灵（实测脚注编号会一路涨到 170+）。
% 钉成 -1：\chapter 明确不步进计数器、也不显示编号，章号完全由 build_pdf.py 显式设定。
\setcounter{secnumdepth}{-1}
\renewcommand{\chaptermark}[1]{\markboth{#1}{}}
\fancypagestyle{plain}{\fancyhf{}\fancyhead[LE,RO]{\thepage}\renewcommand{\headrulewidth}{0pt}}
% 保住中文与西文之间的空格（xeCJK 默认会吃掉）
\xeCJKsetup{CJKspace=true}

% ============ 段落：首行缩进两格 ============
% 中文图书惯例：正文每段首行缩进两格，且**标题后的第一段也缩进**——
% 后者必须显式加载 indentfirst（LaTeX 默认不缩进章标题后的第一段）。
\usepackage{indentfirst}
\usepackage{etoolbox}
\setlength{\parindent}{2\ccwd}
% 以下属于附属文本，一律不缩进（否则会与居中/右对齐/盒子的版式打架）：
\AtBeginEnvironment{quote}{\setlength{\parindent}{0pt}}      % 引文块（含韵文，缩进会破坏断行）
\AtBeginEnvironment{displayeq}{\setlength{\parindent}{0pt}}  % 公式块
\AtBeginEnvironment{infobox}{\setlength{\parindent}{0pt}}    % 知识框（tcolorbox）

% ============ 原版页码（页边灰字） ============
% build_pdf.py 把 <!--p123--> 转成 \origpage{123}；--no-page-marks 时不注入。
\usepackage{xcolor}
% \ifinner 兜底：万一标记落在 minipage/盒子内部，静默跳过而不是中断构建
\newcommand{\origpage}[1]{%
  \ifinner\else\marginpar{\raggedright\footnotesize\sffamily\color{gray}#1}\fi}

% ============ 章节分隔符 • • • ============
\newcommand{\secbreak}{%
  \par\vspace{1.2em}%
  {\centering\large\textbullet\quad\textbullet\quad\textbullet\par}%
  \vspace{1.2em}}

% ============ 推荐语署名（右对齐） ============
% 名家推荐里「——某某，《某某》作者」这类署名行，按中文惯例右对齐。
% 不缩进、也不用 flushright 列表（省掉列表的额外上下间距，单段右对齐即可）。
% 由 tools/infobox.lua 把 ::: {.attribution} 转成本环境。
\newenvironment{attribution}{\par\begingroup\raggedleft}{\par\endgroup}

% ============ 知识框（infobox） ============
% pandoc 默认会丢掉 ::: {.infobox} 这个 Div，靠 tools/infobox.lua 转成此环境。
\usepackage[most]{tcolorbox}
\newenvironment{infobox}{%
  \begin{tcolorbox}[
    enhanced, breakable,
    colback=black!3, colframe=black!25, boxrule=0.4pt, arc=2pt,
    left=6pt, right=6pt, top=5pt, bottom=5pt]%
}{\end{tcolorbox}}

% ============ 图表 ============
% pandoc 默认按原始尺寸插图，宽图会顶出纸面被裁掉（参考项目实测切过 8 行代码）。
\makeatletter
\def\maxwidth{\ifdim\Gin@nat@width>\linewidth\linewidth\else\Gin@nat@width\fi}
\makeatother
\setkeys{Gin}{width=\maxwidth,keepaspectratio}
\usepackage{caption}
% 图号改为「显式编号」：编号写死在图注文字里（由译稿决定），LaTeX 不再自动加「图 N:」。
% 原书的图号体系不是「每张图顺次递增」——它有三类：
%   ① 正常编号图 FIGURE 2.1 …
%   ② **不带编号**的照片/图版（pNNN.jpg，9 张）——自动编号会给它们凭空发号，
%      并把其后所有图号整体推后一位（实测第 4/6/9/10/13 章全部错位）；
%   ③ **同号分幅** FIGURE 2.3A / 2.3B（5 组）——自动编号会把它们拆成两张连号图。
% 两类偏差都会让正文引用的「图 N.M」指错图。改用显式编号后，编号与原文一一对应，
% 且可被 tools/check_figures.py 逐条核对。
% labelsep=none 必须写：光设 labelformat=empty 还会留下一个「: 」分隔符。
\captionsetup{font=small, labelformat=empty, labelsep=none}
\makeatletter
\@addtoreset{figure}{chapter}
\@addtoreset{table}{chapter}
% 脚注也要按章重新编号：原书与书末注释汇编都是每章 1..N，
% 不重置的话正文会出现「170」这种全局编号，与汇编对不上
\@addtoreset{footnote}{chapter}
\makeatother

% ============ 强调与省略号 ============
% 中文加粗改用楷体（Kaiti SC 有真实 Bold 字重，无需合成加粗）。
\newCJKfontfamily\kaiemph{Kaiti SC}[BoldFont={Kaiti SC Bold}]
\renewcommand{\textbf}[1]{{\kaiemph\bfseries #1}}
% 宋体一族会把 U+2026（…）反查成 U+22EF（⋯），复制出来不是同一个码位；
% STFangsong 映射正确。必须用 newCJKfontfamily，否则会被 xeCJK 的标点处理绕过去。
\newCJKfontfamily\ellipsisfont{STFangsong}
\usepackage{newunicodechar}
\newunicodechar{…}{{\ellipsisfont …}}

% ============ 杂项 ============
% 注：正文里出现的希腊字母（如极坐标 (r, θ)）由 build_pdf.py 在数学环境之外
% 替换成 $\theta$。不能用 newunicodechar 处理——它定义的是活动字符，
% 会和 ctex 加载的 microtype 冲突（实测报 \MT@is@char has an extra }）。
% 长 URL 会在脚注里冲出纸面（实测 x 坐标到 601pt > 页宽 595pt，会被裁掉），
% xurl 允许 \url{} 在任意字符处断行
\usepackage{xurl}
% 非 MathML 的行间公式（原书用 <i>/<sub>/<sup> 排的）：居中显示，
% 否则会被 pandoc 排成左对齐的普通段落，与原书不符
\newenvironment{displayeq}{\par\addvspace{0.4em}\centering}{\par\addvspace{0.4em}}
\usepackage{booktabs}
% 目录：紧凑单页、章条目统一样式，并用点引导线连接页码。
% 本书的各级标题在目录中都作为 chapter 条目；ctexbook 默认章条目间距较大，
% 末尾“致谢/注释”会被挤到第二页。压缩章前距但保留清晰行距。
\usepackage{tocloft}
\setcounter{tocdepth}{1}
\setlength{\cftbeforechapskip}{2pt}
\renewcommand{\cftchapleader}{\cftdotfill{\cftdotsep}}
\renewcommand{\cftchapfont}{\normalfont}
\renewcommand{\cftchappagefont}{\normalfont}
% 避免图表浮到下一章
\usepackage{placeins}
